\section{The period quotient and adjoint bigness}\label{sec:quotient}\label{altord:sec:period}
We prove Theorem~\ref{thm:adjoint}. The full period image supplies
an algebraic base, and Section~\ref{sec:rankloss} controls the divisors
introduced by compactification and finite descent. We first show
that those numerical rank drops suffice to remove the divisors
from an adjoint class. This is a section-growth calculation; no
abundance statement for the highest line is needed.

\subsection{Section estimates}\label{altord:sec:estimates}
\begin{lemma}\label{altord:e:volume-lemma}
Let \(T\) be a smooth projective \(t\)-fold, \(A\) a rational divisor,
and \(J\) a nef rational divisor
with \(r=\nu(J)>0\), and \(D=\sum_i d_iE_i\) an effective integral divisor
with smooth prime components. Suppose
\[
 \nu(J|_{E_i})<r\quad\text{for every }i.
\]
If \(A+D\) is big, then \(A+jJ\) is big for every sufficiently large integer
\(j\).
\end{lemma}

\begin{proof}
Choose an ample rational divisor \(H_1\) such that \(A+D-H_1\) is
effective. In this proof \(m\) tends to infinity through sufficiently
divisible positive integers. For each fixed \(j\geq0\), asymptotic Riemann--Roch for the ample
divisor \(H_1+jJ\) gives
\[
 \liminf_{m\to\infty}
 \frac{t!}{m^t}h^0\bigl(T,m(A+D+jJ)\bigr)
 \ \geq\ (H_1+jJ)^t\ \geq c j^r
\]
for a constant \(c>0\) and all \(j\geq1\).
The last inequality follows from \(J^rH_1^{t-r}>0\).

Remove \(mD\) one prime divisor at a time. Each quotient in a divisor
exact sequence is a line bundle on some \(E_i\), of the form
\[
 \mathcal O_{E_i}
 \left(m(A+D+jJ)-\sum_h k_hE_h\right),
 \qquad 0\leq k_h\leq md_h.
\]
There is a single ample rational divisor \(H_2\) on \(T\) for which every
such quotient has at most
\(h^0(E_i,m(H_2+jJ)|_{E_i})\) sections, in divisible degrees and for every
\(j\geq0\). To see this, choose an ample Cartier divisor \(B_0\) with a
globally generated multiple clearing the denominator of
\(B_0-(A+D)\), and ample Cartier divisors \(B_h\) such that both \(B_h\)
and \(B_h+E_h\) are globally generated. Enlarge \(B_0\) if necessary and
take \(H_2=B_0+\sum_h d_hB_h\). The difference, restricted to \(E_i\), is
\[
 m(B_0-A-D)+
 \sum_h\bigl((md_h-k_h)B_h+k_h(B_h+E_h)\bigr).
\]
It has a section nonzero on \(E_i\), by global generation and
multiplication, giving the claimed bound.

There are \(m\sum_i d_i\) quotient terms. For fixed \(j\),
asymptotic Riemann--Roch on each \(E_i\) bounds their total contribution on
the scale \(m^t/t!\) by
\[
 t\sum_i d_i\,(H_2+jJ)^{t-1}\cdot E_i.
\]
This is \(O(j^{r-1})\), since every intersection containing at least
\(r\) factors of \(J|_{E_i}\) vanishes. Consequently
\[
 \liminf_{m\to\infty}
 \frac{t!}{m^t}h^0\bigl(T,m(A+jJ)\bigr)
 \ \geq\ c j^r-O(j^{r-1})>0
\]
for all sufficiently large \(j\), proving bigness.
\end{proof}

\begin{lemma}\label{altord:e:exceptional-pullback}
Let \(\pi:T\to S\) be a surjective generically finite morphism of smooth
projective varieties. Let \(E\geq0\) be a rational divisor on \(T\) whose
components have image of codimension at least two in \(S\). For a rational
divisor \(D\) on \(S\), the divisor \(\pi^*D+E\) is big if and only if
\(D\) is big.
\end{lemma}

\begin{proof}
Factor \(\pi\) as \(T\xrightarrow{\rho}\bar S\xrightarrow{\tau}S\),
where \(\bar S\) is the normalization of \(S\) in \(\C(T)\).
Then \(\rho\) is birational and \(\tau\) is finite. Every component of
\(E\) is \(\rho\)-exceptional, since a finite map preserves dimensions.
Normality gives
\(\rho_*\mathcal O_T(mE)=\mathcal O_{\bar S}\) in divisible degrees:
a rational function with poles only on exceptional divisors has no
codimension-one pole downstairs. Projection formula identifies the
sections of \(m(\pi^*D+E)\) with those of \(m\tau^*D\).
Finally, bigness is invariant under surjective finite pullback.
\end{proof}


\subsection{Construction and descent}
\begin{proof}[Proof of Theorem~\ref{altord:p:observation}]
We may first make the birational modification allowed in the hypothesis.
Replace $\mathbb V$ by the minimal rational subvariation from
Lemma~\ref{altord:h:minimal}, and equip it with the intersection lattice and
the restricted polarization.  Its deepest Hodge line and the line
$J_{\widehat Y}$ are unchanged.  We continue to denote this variation
by $\mathbb V$.

\medskip\noindent
\textbf{The actual period image and its finite quotient.}
Choose a rational reference space with lattice and polarization, and
let $\mathcal D$ be the connected period-domain component containing a
chosen lift of the periods. Let $\Gamma$ be the arithmetic group of
integral polarizing isometries preserving this component; it contains
the monodromy of $\mathbb V$. Let $\Gamma'\triangleleft\Gamma$ be a
normal neat subgroup of finite index.  The kernel of
\[
  \pi_1(U)\longrightarrow\Gamma/\Gamma'
\]
defines a connected finite \'etale Galois cover $U'\to U$.  Denote its
deck group by $H$; thus $H$ is the image of this homomorphism, not
necessarily all of $\Gamma/\Gamma'$.

The period map of $\mathbb V|_{U'}$ factors algebraically through the
normalization $B$ of the closure of its image in
$\Gamma'\backslash\mathcal D$.  The space $B$ is a normal
quasiprojective variety, and on a smooth dense open subset its period
realization has generically immersive differential; this is the
period-image theorem \cite[Theorem~1.1]{BakkerBrunebarbeTsimerman}.  Here and below the function
field of $B$ is the function field of this normalized \emph{actual
image}.  We do not replace it by its relative algebraic closure in
$\C(U')$.

Neatness makes the arithmetic action free.  The universal local system
on the quotient is
\[
  (\mathcal D\times V_{\Z})/\Gamma',
\]
with its tautological filtration. On the smooth open of $B$ this
filtration is horizontal: the dominant map $U'\to B$ is generically
submersive and its pullback is the original variation, so Griffiths
transversality holds on a dense open and hence everywhere. Thus the
restricted local system and filtration form a polarizable integral
variation. Moreover, the action of
$\Gamma/\Gamma'$ on this local system is an honest action:
\[
  [x,v]\longmapsto[\gamma x,\gamma v].
\]
Normality of $\Gamma'$ makes the formula independent of the
representative $\gamma$.  Restricting gives an action of $H$ on the
variation over $B$, covering its action on $B$. The action on $B$ is
algebraic: apply uniqueness in \cite[Theorem~1.1]{BakkerBrunebarbeTsimerman} to the period map
and its expression obtained by a deck transformation on $U'$ and the
corresponding arithmetic translation on the target, and then lift to
the normalization. Let
\[
  G_f=H/\ker(H\longrightarrow\operatorname{Aut}(B))
\]
be the effective group on the image.  The full group $H$ will still be
retained for its action on the variation.

Choose a $G_f$-equivariant smooth projective compactification $S'$ of a
smooth dense invariant open of $B$, with snc boundary, by equivariant
compactification and resolution.  All boundary monodromies of its
variation are unipotent.  Indeed they are quasi-unipotent by the
monodromy theorem, and they belong to the neat group $\Gamma'$; their
root-of-unity eigenvalues must therefore all be one.  Let $J'$ denote
the extended deepest Hodge line and put $L'=cJ'$.  This is nef by
Theorem~\ref{altord:h:inputs}.  The $H$-action extends to $J'$ by
functoriality of the unipotent Deligne and Hodge extensions.

Let $q:S'\to Q_0=S'/G_f$ be the finite quotient.  Although an element
of $\ker(H\to G_f)$ may act nontrivially on $J'$, it acts on its fibres
through finite-order characters.  Taking a tensor power divisible by
$|H|$ kills the fibre action of every stabilizer in $H$, including this
kernel.  That tensor power therefore descends to an actual line bundle
on $Q_0$ by finite-group descent.  Consequently there is a rational
line bundle $\overline L$ on $Q_0$ with
\[
  q^*\overline L\sim_{\Q}L'.
\]
It is nef, since its pullback by the finite surjective map $q$ is nef.
Choose a smooth projective resolution $s:S_0\to Q_0$ and put
$L_0=s^*\overline L$.

The equivariant period map on $U'$ descends to a dominant rational map
$Y\dashrightarrow Q_0$, and hence to $S_0$.  Resolve its graph by a
smooth projective birational modification of $Y$, obtaining
\[
  p_0:Y\longrightarrow S_0.
\]
We claim that
\begin{equation}
  M\sim_{\Q}p_0^*L_0.
  \label{altord:p:line-descent}
\end{equation}
To verify this claim, choose a common smooth projective alteration $W$
dominating the alteration in the hypothesis and the level cover, and
resolve the graph of its rational map to $S'$.  Write
$\mu:W\to Y$ and $t:W\to S'$ for the resulting morphisms.  Arrange
snc boundaries by further resolution, containing the inverse images of
the boundaries in question.  The pulled-back variations agree on a
dense open subset.  Unipotent Deligne extensions and their extended
Hodge filtrations commute with these pullbacks between smooth snc
pairs; see \cite[Lemma~5.6]{BFMT}, or the unipotent extension theory
\cite{Schmid,CKS}.  The extended top line $J_W$ is consequently the
pullback of the top line on either model.  The hypothesis and the
commutative maps to $Q_0$ give
\[
  \mu^*M\sim_{\Q}cJ_W
  \sim_{\Q}t^*L'
  \sim_{\Q}\mu^*p_0^*L_0.
\]
Pullback by a proper generically finite map is injective on rational
line-bundle classes: applying divisor pushforward, or the norm of a
principal rational function, multiplies the original class by the
degree.  This proves \eqref{altord:p:line-descent}.  In particular no boundary
twist has been inferred from an equality merely on the open locus.

If $\dim S_0=0$, equation~\eqref{altord:p:line-descent} gives
$M\sim_{\Q}0$, and the theorem follows.  Assume henceforth that
$\dim S_0>0$.

\medskip\noindent
\textbf{A common equivariant model and log general type.}
Normalize $S_0$ in $\C(S')$, and take a $G_f$-equivariant smooth
projective resolution of the resulting model and its birational map
to $S'$.  Denote the resulting maps by
\[
  \eta:T\longrightarrow S',\qquad
  \pi:T\longrightarrow S_0.
\]
Thus $\eta$ is birational, $\pi$ is generically finite with Galois
function-field group $G_f$, and $G_f$ acts on $T$ over $S_0$.
Choose the resolution so that there is a $G_f$-invariant snc boundary
$D_T$ containing the complement of the open period locus and the
support of the Jacobian divisor of $\pi$.  The pulled-back variation
has unipotent boundary monodromy.  If $J_T$ is its extended top line,
then
\begin{equation}
  L_T:=cJ_T\sim_{\Q}\eta^*L'
     \sim_{\Q}\pi^*L_0.
  \label{altord:p:T-line}
\end{equation}

Let $D_N$ be the reduced sum of those components of $D_T$ having
nonzero monodromy logarithm.  The variation extends as a pure polarized
variation over
\[
  U_N=T\setminus D_N.
\]
Indeed, near a point where all the remaining logarithms are zero, the
nilpotent-orbit extension is a pure polarized Hodge filtration; the
several-variable extension gives the same statement at crossings
\cite{Schmid,CKS}.  These extensions are unique and preserve the
$H$-equivariance of the variation.  The divisor $D_N$ is snc, since it
is a union of components of $D_T$.

The full period map on $U_N$ has generically immersive differential,
because $T$ is birational to $B$.  The log-general-type theorem for a
polarized variation with generically immersive period map therefore
gives
\begin{equation}
  K_T+D_N\quad\text{big}.
  \label{altord:p:log-big}
\end{equation}
Here the applicable theorem is \cite[Theorem~1.1]{BrunebarbeCadorel}.  Its boundary is precisely
$D_N$, because the pure variation has just been extended over all the
other boundary components.

We now have precisely the geometric data of
Theorem~\ref{thm:rankloss}: $T$ is birational to the actual normalized
period image, the full period map is generically immersive, and the
retained finite group $H$ acts on its integral polarized variation.
The effective group on $T$ is $G_f$. Minimality is preserved here:
the variation on $T$ and the original variation have the same
Hodge structures at Hodge-generic points, since the original source
dominates the period image. Consequently
\begin{equation}
 r:=\nu(L_T)>0,
 \qquad \nu(L_T|_E)<r
 \label{altord:p:rank-drop}
\end{equation}
for every component of $D_N$ and every divisorial ramification
component of $\pi$ mapping onto a divisor of $S_0$. For the latter,
the inertia group of the Galois extension $\C(T)/\C(S_0)$ contains
a nonidentity element of $G_f$ fixing the divisor pointwise. These
fixed loci were included in $D_T$. Positive rational rescaling from
$J_T$ to $L_T$ leaves all these numerical dimensions unchanged.
The use of the actual image, before Stein factorization, is exactly
the function-field hypothesis needed in the finite-inertia part of
Theorem~\ref{thm:rankloss}.

\medskip\noindent
\textbf{Removal of the boundary and ramification.}
Write the effective Jacobian divisor as
\[
  R_\pi=K_T-\pi^*K_{S_0}
       =R_{\mathrm{div}}+R_{\mathrm{exc}},
\]
where $R_{\mathrm{div}}$ is supported on components mapping onto
divisors of $S_0$, and each component of $R_{\mathrm{exc}}$ has image
of codimension at least two.  In the volume
Lemma~\ref{altord:e:volume-lemma}, take
\[
  J=L_T,\qquad
  D=D_N+R_{\mathrm{div}},\qquad
  A=K_T-R_{\mathrm{div}}
   =\pi^*K_{S_0}+R_{\mathrm{exc}}.
\]
All components of $D$ satisfy the rank inequality above, and
$A+D=K_T+D_N$ is big by \eqref{altord:p:log-big}.  The volume lemma gives
\[
  A+jL_T
  \sim_{\Q}\pi^*(K_{S_0}+jL_0)+R_{\mathrm{exc}}
  \quad\text{big for }j\gg0.
\]
Lemma~\ref{altord:e:exceptional-pullback} removes the effective divisors
exceptional over $S_0$ and descends bigness through the generically
finite map.  Therefore
\begin{equation}
  K_{S_0}+jL_0\quad\text{is big for every sufficiently large integer }j.
  \label{altord:p:adjoint-big}
\end{equation}

\medskip\noindent
\textbf{Connected fibres.}
Finally take the Stein factorization of $p_0$, and resolve its finite
normal intermediate variety to a smooth projective variety $S$.
Resolving the induced rational map from $Y$ gives, by further smooth
projective birational modification, a morphism $p:Y\to S$ with
connected fibres. Indeed, its Stein factor is finite birational over
the normal variety $S$, because $\C(S)$ is already algebraically closed
in $\C(Y)$, and is therefore $S$ itself.
If $\sigma:S\to S_0$ is the resulting generically
finite morphism, set $L=\sigma^*L_0$.  It is nef, and
\eqref{altord:p:line-descent} pulls back to $M\sim_{\Q}p^*L$.
Since both $S$ and $S_0$ are smooth, the Jacobian divisor of $\sigma$
is effective.  Thus
\[
  K_S+jL
  =\sigma^*(K_{S_0}+jL_0)+R_\sigma
\]
is big for $j\gg0$ by \eqref{altord:p:adjoint-big} and preservation of
bigness under generically finite pullback.  This proves the theorem.
\end{proof}
